\scanpage{516}
\mspage{427}
% B01
\uline{Dém} \inserted{de 4.15.1 et 4.15.2.} Il est clair que \deleted{\unsure{8}}4.15.1 est contenu dans 4.15.2 a) et c). Nous allons donc prouver 4.15.2.

% B02
a) Soit $(x_i)_{i\in I}$ un système inductif dans $E_s$, avec $I$ grd devant $\pi_1$, et $x$ sa limite ind. dans $E_s$ \inserted{(qui existe, car $E_s$ satisfait à $L_{\pi_\lambda}$ donc \unsure{à|:} $L_{\pi_0}$)}. Il faut prouver que c'est une limite dans $E$, donc que pour tt $y\in\mathrm{Ob}\,E$, la flèche canonique
% B03
{\csname tagsleft@true\endcsname
\[\deleted{\unsure{\lim}}\,\mathrm{Hom}_E(\unsure{x|y},y)\unsure{\relbar|\longrightarrow}\varprojlim_I\mathrm{Hom}_E(x_i,y)\tag*{\unsure{(a)|(*)}}\]}
% B04
est un iso. Or $\mathrm{Hom}_E(x,y)$ et les $\mathrm{Hom}_E(x_i,y)$ sont des objets de $\mathrm{Ens}_{/\mathrm{Hom}_C(s,t)}$, où $s=f(x)$, $t=f(y)$, et les flèches de transition sont des flèches de $\mathrm{Ens}_{/\mathrm{Hom}_C(s,t)}$. Cela montre que la bijectivité de \unsure{(a)|(*)} revient à \unsure{celles|celle} des
% B05
{\csname tagsleft@true\endcsname
\[\mathrm{Hom}_\lambda(x,y)\relbar\varprojlim_I\mathrm{Hom}_\lambda(x_i,y)\tag*{\unsure{(a)|(* )}}\]}
% B06
pour tout $\lambda\in\mathrm{Hom}_C(x,y)$. Pour $\lambda$ fixé, et $y$ fixé dans $\mathrm{Ob}\,E_t$, la bijectivité en \unsure{qs|qn}, pour tout système inductif grd devant $\pi_1$ dans $E_s$, de limite $x$, signifie que le foncteur
% B07
\[x\longmapsto\mathrm{Hom}_\lambda(x,y):\quad\deleted{\unsure{E| \mathrm{Ens}}}E_s\relbar(\mathrm{Ens})^\circ\]
% B08
est $\pi_{\unsure{0|1}}$-accessible. Or c'est bien le cas, par hypothèse sur $\pi_1$, quel que soit la flèche \deleted{\unsure{$\lambda$ de sou}} $\lambda:s\to t$ de source $s$, et quel que soit $y\in\mathrm{Ob}\,E_t$.

% archive sheet number: 515
